% !TEX program = lualatex
\documentclass[11pt,a4paper]{ltjsarticle}
\usepackage{amsmath,amssymb,amsthm,mathtools,luatexja-fontspec}
\usepackage[margin=25mm]{geometry}
\usepackage[colorlinks=true,urlcolor=blue,linkcolor=blue]{hyperref}
\setmainjfont{HaranoAjiMincho-Regular}
\newtheorem{theorem}{定理}[section]
\newtheorem{proposition}[theorem]{命題}
\newtheorem{lemma}[theorem]{補題}
\newtheorem{corollary}[theorem]{系}
\newtheorem{remark}[theorem]{注意}
\newcommand{\CC}{\mathbb C}
\newcommand{\PP}{\mathbb P}
\newcommand{\Sym}{\operatorname{Sym}}
\newcommand{\Spec}{\operatorname{Spec}}
\newcommand{\GL}{\operatorname{GL}}
\renewcommand{\proofname}{証明}
\title{標識付きHesse平面切断の臨界幾何}
\date{2026年11月13日}
\begin{document}
\maketitle
\begin{abstract}
三次三様体の平面切断を、三変数三次式の標識付きHesse切片を通して調べる。
臨界平面の九係数模型から六次Hessian因子を得て、その絶対既約性、正規性、有理性を証明する。
核上三次の消失集合は四つの有理的成分を持つ。
孤立A2、A3点で十次元の三次式変形が与える局所unfoldingと、
臨界値像の正規化・導手を具体的に計算する。
幾何は形式的または指定点での完備局所的なものであり、有限量子係数の圏同値を前提にしない。
\end{abstract}
\section{平面関数と臨界条件}
標数零で考える。$D=x_1^3+x_2^3+x_3^3$とし、巡回する$i,j,k$に対して
\begin{equation}\label{mh:family}
 F_{s,P,Q,T}=D/3+u^3+v^3+s x_1x_2x_3+
 \sum_i(P_i u+Q_i v)x_jx_k+\sum_i T_i x_iuv
\end{equation}
を取る。$u=\ell=\sum\lambda_i x_i$, $v=m=\sum\mu_i x_i$は平面の局所座標である。
原点に合同な三変数三次式$f$には、一意な$N\equiv I$と$\sigma$が存在して
$f(Nx)=D/3+\sigma x_1x_2x_3$となる。
実際、九つの行列方向の微分は$x_i^2x_j$の九単項式であり、残る$xyz$方向と合わせて全三次式の基底をなす。
形式逆関数定理を用いれば、全Artin基底に両立する一意な規格化を得る。
この規格化を制限式へ適用した係数を$\sigma(s,P,Q,T,\lambda,\mu)$と書く。
$\sigma(s,P,Q,T,0,0)=s$であり、全変数の通常次数で
\begin{equation}\label{mh:jet}
 \sigma=s+\sum_i(P_i\lambda_i+Q_i\mu_i)
 +6(\lambda_1\lambda_2\lambda_3+\mu_1\mu_2\mu_3)
 +\sum_{\{i,j,k\}=\{1,2,3\}}T_k\lambda_i\mu_j+O(4).
\end{equation}
二次の規格化は非対角行列であり、その$sxyz$への作用に$xyz$項はないので、この三次jetは規格化後にも成立する。

\begin{proposition}[内在的臨界条件]\label{mh:intrinsic}
滑らかな平面三次曲線$E=V(f)\subset W\subset X=V(F)$に対し、平面移動による曲線の変形が零となる条件は
\[
 F_u|_W,F_v|_W\in (\partial f)_2.
\]
これは極化写像$V\to\Sym^2 W^*$の階数が三であること、
および$0\to T_E\to T_X|_E\to N_{E/X}\simeq\mathcal O_E(1)^2\to0$の分裂と同値である。
このとき$V=W\oplus\ker DF_W$という標準補空間を取り、
\[
 F=f(x)+\sum_i x_iq_i(u,v)+g(u,v)
\]
と書ける。特に平方線形混合項を含まない。
\end{proposition}
\begin{proof}
三つの偏微分は二次式の正則列であり、Jacobian環のHilbert関数は$(1,3,3,1)$である。
その完全対合$R_1\otimes R_2\to R_3$から、任意の一次式$a,b$に対し
$aF_u+bF_v$がJacobianイデアルに入る条件は、表示の二つの二次類の消滅と等しい。
これは平面切断のKodaira--Spencer写像の消滅である。
同じ写像は法束列の連結写像であり、$T_E\simeq\mathcal O_E$とSerre双対を用いると、
その消滅は拡張類の消滅に等しい。三つの$x$偏微分は独立なので極化の階数は少なくとも三であり、
残る二偏微分がその像に入ることが階数三と等しい。
核は$W$と交わらず、核方向を$u,v$に選べば$x^2u,x^2v$項が消える。
\end{proof}

以後、臨界点での三次式を
\begin{equation}\label{mh:criticalfamily}
 F_Q=D/3+u^3+v^3+\sum_i x_i(a_i u^2+d_iuv+c_i v^2),\qquad
 Q_i=\begin{pmatrix}2a_i&d_i\\d_i&2c_i\end{pmatrix}
\end{equation}
とする。平面変数を交互順序$n=(\lambda_1,\mu_1,\lambda_2,\mu_2,\lambda_3,\mu_3)$に並べると、
$\sigma$のHessianおよび三次項は
\begin{equation}\label{mh:hessian}
 H=\begin{pmatrix}0&Q_3&Q_2\\Q_3&0&Q_1\\Q_2&Q_1&0\end{pmatrix},\qquad
 \Gamma(n)=6(\lambda_1\lambda_2\lambda_3+\mu_1\mu_2\mu_3).
\end{equation}
入力の平面次数が二以上なので、Hesse規格化の非線形効果は四次以上であり、これは正確な二、三次項である。

\section{六次Hessian因子}
\begin{theorem}\label{mh:delta}
九係数空間で$\Delta=\det H$と置くと
\[
 \Delta=4\prod_i(4a_ic_i-d_i^2)-16\det(a,d,c)^2.
\]
$\mathcal D=V(\Delta)\subset\mathbb A^9$は絶対既約な正規有理八次元超曲面で、
特異集合の次元は六である。原点での完備局所環も正規整域である。
階数五で$\ker H=\langle(n_1,n_2,n_3)\rangle$, $n_i\in K^2$と書けば、
$\mathcal D$が特異である条件は非零な$n_i$が高々一つであることである。
\end{theorem}
\begin{proof}
対称二行行列$A,B,C$について、直接Schur補元を取ると
\[
 \det\begin{pmatrix}0&C&B\\C&0&A\\B&A&0\end{pmatrix}
 =4\det A\det B\det C-
 \det\begin{pmatrix}A_{11}&A_{12}&A_{22}\\B_{11}&B_{12}&B_{22}\\C_{11}&C_{12}&C_{22}\end{pmatrix}^{\!2}.
\]
$C$可逆の稠密開集合で合同変換により$C=I$とし、さらに$A$を対角化すると、
右辺は$\det(AB+BA)$の計算に帰着する。多項式恒等式なので全体に延長する。
$Q_i$を代入すると表示式を得る。

$A=\left(\begin{smallmatrix}x&y\\y&z\end{smallmatrix}\right)$,
$B=\left(\begin{smallmatrix}u&v\\v&w\end{smallmatrix}\right)$,
$C=\left(\begin{smallmatrix}r&s\\s&t\end{smallmatrix}\right)$とする。
上式を$R=\CC[x,y,z,u,v,w]$上の$r,s,t$に関する二次式と見る。
$r^2,t^2$係数はそれぞれ$-(yw-zv)^2,-(xv-yu)^2$であり、二つの一次原始既約式は非同伴なので二次式は原始的である。
$A=I,B=\operatorname{diag}(1,2)$で二次式は$8rt-9s^2$となり、その行列式は$144$である。
従って$\operatorname{Frac}R$の代数閉包上でも階数三で、一次式二つの積にはならない。
Gaussの補題から絶対既約性を得る。

$A$可逆の開集合では$\operatorname{rank}H=4+\operatorname{rank}(CA^{-1}B+BA^{-1}C)$。
$A=I$に帰着し、$B=\left(\begin{smallmatrix}x&y\\y&z\end{smallmatrix}\right)$に対する
$C\mapsto BC+CB$の行列は
\[
 \begin{pmatrix}2x&2y&0\\y&x+z&y\\0&2y&2z\end{pmatrix},\qquad
 \det=4(xz-y^2)(x+z).
\]
その階数が一以下なら二次小行列式から$B=0$となる。
よって階数三、二、零の層で$H$の階数四以下の集合の次元はいずれも高々六である。
三行列とも特異な補集合も次元高々六である。
階数五では$\operatorname{adj}H=\kappa kk^t$なので、$Q_i$方向の微分は
$2\kappa n_j^t(\delta Q_i)n_k$である。
任意の対称行列$S$について$e^tSf=0$となるのは$e=0$または$f=0$の場合だけであり、特異条件を得る。
例えば$n_1\ne0,n_2=n_3=0$なら$Q_2n_1=Q_3n_1=0$であり、
核直線の一次元、$Q_1$の三次元、$Q_2,Q_3$それぞれの一次元を足して六次元以下である。
一方$Q_1=0$の六次元空間は階数四以下なので全体が特異である。
従って特異集合の次元は正確に六であり、整超曲面の$(S_2)$と余次元二以上の特異集合から$(R_1)$を得て正規性が従う。
$\Delta$は斉次既約なので完備環の随伴次数環も$\CC[Q]/(\Delta)$という整域である。
分離的フィルトレーションから完備環も整域であり、excellent正規局所環の完備化の正規性を用いる。

有理性には核incidenceを用いる。$[k]\in\PP^5$に対する
$L_k:\Sym^2(K^2)^{\oplus3}\to K^6$, $Q\mapsto H(Q)k$の全射開集合上で、
incidenceは階数三のベクトル束である。
\[
 k=(1,0,0,1,1,1)^t,\quad
 Q_1=\begin{pmatrix}1&1\\1&-1\end{pmatrix},\quad
 Q_2=\begin{pmatrix}-1&1\\1&2\end{pmatrix},\quad
 Q_3=\begin{pmatrix}-2&0\\0&-3\end{pmatrix}
\]
では$L_k$の列$1,2,3,4,5,8$の行列式は$2$、$H$の第一行列を除く小行列式は$-12$である。
階数五の上では核直線が一意なので、この有理八次元束は$\mathcal D$と双有理である。
\end{proof}

\section{核上三次の四成分}
この節の行列成分は$Q_i=\left(\begin{smallmatrix}A_i&B_i\\B_i&C_i\end{smallmatrix}\right)$と記す。
$\mathcal D_5=\{\Delta=0,\operatorname{rank}H=5\}$上で
$\Gamma|_{\ker H}=0$の零スキームを$\mathcal C_5$、そのスキーム論的閉包を$\overline{\mathcal C}_5$とする。
\begin{theorem}\label{mh:four}
$\overline{\mathcal C}_5$は被約純七次元で、ちょうど四つの有理的既約成分
$M,E_1,E_2,E_3$を持つ。$T=\det(A,B,C)$と置くと
\[
 E_i=V(\det Q_i,T),\qquad
 \deg\PP(M)=27,\quad\deg\PP(E_i)=6.
\]
$M$は三ブロック全て非零な核を持つ部分の閉包である。
随伴行列の第$j$列を$v_j$とすると全閉包の有限イデアル表示は
\[
 (\Delta,\Gamma(v_1),\ldots,\Gamma(v_6)):I_5(H)^\infty.
\]
\end{theorem}
\begin{proof}
核incidenceの階数をまず全て数える。非零ブロックが三つなら$\operatorname{rank}L_k=6$、
二つで独立なら五、二つで比例なら四、一つなら四であり、対応する射影基底の次元は$5,3,2,1$である。
例えば独立な二ブロックを$e_1,e_2$、第三を$(r,s)$にした六方程式では、六次小行列式に$-2r,-2s$が現れる。
三つとも比例する場合は$n_i=t_ie_1$とすれば二組の三次行列式が$2t_1t_2t_3$である。
ただしこのとき全ての$Q_i e_1=0$で$H$の階数は三以下である。
二つで独立かつ第三零なら$Q_3=0$となり$H$の階数は四以下である。
残る場合の核条件は$Q_3n=0,t_1Q_2n+t_2Q_1n=0$、
または$Q_2n_1=Q_3n_1=0$であり、表示の階数が従う。

$X=V(\lambda_1\lambda_2\lambda_3+\mu_1\mu_2\mu_3)\subset\PP^5$は原始一次式として既約であり、
$\lambda_1\lambda_2\ne0$上で有理的である。特異集合は$\lambda_i$座標点と$\mu_j$座標点を結ぶ九本の直線である。
従って正規で、三ブロック非零の開集合$U$上では滑らかである。
$X\cap U$上の階数三束は七次元の有理的主成分を与える。
核の一ブロックが零で他の二つが比例する各場合は、二次元の有理的基底上の階数五束を与え、七次元成分となる。
一ブロックのみ非零の部分は六次元以下である。
$\mathcal D_5$はCohen--Macaulayで、$\Gamma$は零でない切断なのでその零スキームは純七次元である。
従って以上で成分を尽くす。

例えば$Q_3$階数一なら合同変換で$Q_3=\operatorname{diag}(0,e)$とでき、
$T=e(A_1B_2-B_1A_2)$である。これから$V(\det Q_3,T)$の$Q_3\ne0$部分は既約で一般点被約、
$Q_3=0$部分は六次元である。二、三次式の完全交差なので埋込み成分はなく、全体が整で第三境界成分に一致する。
他も同様である。
主成分では滑らかな$X\cap U$により一般点の重複度は一である。
境界では
\[
 Q_1=\begin{pmatrix}1&1\\1&2\end{pmatrix},\quad
 Q_2=\begin{pmatrix}-1&-1\\-1&3\end{pmatrix},\quad
 Q_3=\begin{pmatrix}0&0\\0&1\end{pmatrix},\quad k=(1,0,1,0,0,0)
\]
を用いる。第一行列を除く小行列$J$は$\det J=-8$、$(J^{-1})_{\lambda_3,\lambda_3}=-1/8$。
$\delta A_1=1$は$T_Q\mathcal D$に属し、正規化した核の微分は$\delta\lambda_3=1/8$、
従って$d\Gamma=3/4\ne0$である。Cohen--Macaulay性から全体が被約である。
階数五チャート上では随伴行列の列が核の生成元を張るため、飽和イデアルの式も従う。

次数は$\PP^8_Q\times\PP^5_k$上の六本の双次数$(1,1)$方程式$Hk=0$で計算する。
上の階数表からincidenceは純七次元の整完全交差であり、その類は$(A+B)^6$である。
$\Gamma=0$との交差を$Q$側に押すと$3B(A+B)^6$の$B^5$係数は$45A^2$である。
階数四以下に写る部分は像の射影次元が五以下なので寄与しない。
三境界成分はそれぞれ次数六、各成分の写像次数・重複度は一なので、主成分は$45-18=27$となる。
\end{proof}

\begin{proposition}[一ブロック核の非孤立極限]\label{mh:nonisolated}
階数五で核が$(n,0,0)$なら、平面関数は形式右同値で$\sigma=q_5(y)$となる。
従って$h^2=3\sigma$のファイバーは五次元通常二重点と形式直線の積であり、
$T^1=K[[z]]$である。
\end{proposition}
\begin{proof}
$Q_2n=Q_3n=0$である。$u=z\lambda x_1,v=z\mu x_1$への制限は
$c(z)x_1^3+(x_2^3+x_3^3)/3$、$c(0)=1/3$となる。
対角規格化でHesse係数は恒等的に零であり、二つの正規偏微分も$x_1^2$の倍数なので全平面微分が零である。
高速五方向の停留解は一意なのでこの曲線に一致する。形式Morse分裂により残余関数は厳密に零である。
\end{proof}

\section{孤立低次特異点と標的変形}
全ての完備近傍は、指定した三次式と平面の点を中心に取る。
$\rho$を可逆にした後でFermat原点の完備化をそのまま用いるのではない。
\begin{theorem}\label{mh:versal}
\eqref{mh:criticalfamily}が滑らかな三次三様体を与え、$H$が階数五であるとする。
核方向三次が非零なら$A_2$型、三次が零でMorse消去後の四次係数が非零なら$A_3$型である。
いずれも十次元の三次式の標的変形はminiversal基底を含む。
正確には、$h$を固定する相対形式座標と標的座標により、それぞれ
\[
 h^2=3\{q_5(y)+\eta z^3+az+b\},\qquad
 h^2=3\{q_5(y)+\kappa z^4+az^2+bz+c\}
\]
と書ける。ここで$\eta,\kappa\in K^\times$、$q_5$は非退化二次形式であり、
それらの根の添加や二次形式の分裂を仮定しない。
\end{theorem}
\begin{proof}
滑らかな五変数三次式の偏微分は二次正則列なので一次係数のsyzygyはない。
従って二十五の$w_j\partial_{w_i}F$は独立で、その補空間が十次元形式sliceを与える。
相対Morse消去で$\sigma=q_5(y)+g(t,z)$とする。
中央の$g$は非零な三次または四次から始まるので、定数項一の単元の根で高次項を除ける。
標的イデアルの各小拡大で生じる誤差を、三次の場合
$e_0+e_1z+z^2R(z)$、四次の場合$e_0+e_1z+e_2z^2+z^3R(z)$と分ける。
低次項を係数へ吸収し、残りを$z\mapsto z-R/(3\eta)$または$z-R/(4\kappa)$で除く。
標的極大イデアルに属する定数項の平行移動も許し、全局所位相で帰納する。
この操作は方程式の単元倍を必要とせず、$h$を固定する。

係数の微分が独立であることを示す。核ベクトルを$k$とすると、
$xyz$および$ux_jx_k,vx_jx_k$方向の残余Jacobian商での像は、定数項一、
および一次係数$k_{\lambda_i},k_{\mu_i}$である。
Hesse行列が$I+O(n^2)$、高速停留解が$O(z^2)$なので、この主係数は変わらない。
核は非零なので三次の場合は二方向が独立である。
四次の場合、命題\ref{mh:nonisolated}から二つ以上の非零ブロック$n_j,n_k$がある。
二次形式$q(u,v)$を選び、その偏極$B_q(n_j,n_k)\ne0$とできる。
方向$x_iq(u,v)$は純二次の非零類$B_q(n_j,n_k)z^2$を与える。
従って三方向が下三角に独立である。
周囲の座標軌道方向は反対向きの平面移動で相殺され、平面関数の偏微分の線形結合となるので残余Jacobian商で零である。
よって独立な類は十次元sliceに降り、逆関数定理により係数を標的座標の一部にできる。
\end{proof}

\begin{proposition}[非空性と四成分の一般型]\label{mh:examples}
$\mathcal D$の一般点は$A_2$型である。
$\overline{\mathcal C}_5$の四成分はいずれも一般点で$A_3$型である。
\end{proposition}
\begin{proof}
$K=\CC((\rho))$上で、$A_2$の例は
\[
 Q_1=\rho I,\quad Q_2=\rho\operatorname{diag}(1,4),\quad
 Q_3=\rho\begin{pmatrix}5&4\\4&5\end{pmatrix},\quad
 k=(-2,1,2,-4,6,3).
\]
最後の行列を除く五次小行列式は$-90\rho^5$、$\Gamma(k)=-216$である。
$\rho=0$でFermat式に還元されるため$F_Q$は幾何学的に滑らかである。
条件は開で、定理\ref{mh:delta}の既約性から稠密である。
この弧は原点を通るので原点完備環の通常の$\Spec$でも非空稠密開集合を与える。

主成分の$A_3$の例は
\[
 Q_1=\rho I,\quad Q_2=\rho\begin{pmatrix}0&-1\\-1&-1\end{pmatrix},\quad
 Q_3=\rho\begin{pmatrix}-3&1\\1&0\end{pmatrix},\quad
 k=(1,1,1,2,2,-1).
\]
$H=\rho H_0$の最初の五次小行列式は$2\rho^5$、$\Gamma(k)=0$である。
核方向の一次式を$L=x_1+x_2+2x_3$, $M=x_1+2x_2-x_3$とする。
\[
 p_2=\tfrac12\sum_i x_i\{(Q_i/\rho)_{11}L^2+
 2(Q_i/\rho)_{12}LM+(Q_i/\rho)_{22}M^2\}.
\]
$N=I+\rho z^2 B+O(z^3)$による二次規格化は
\[
 B=\begin{pmatrix}-1&-3/2&-1/2\\5/2&4&1/2\\5/2&3/2&8\end{pmatrix},
 \qquad p_2+\sum_i x_i^2(Bx)_i=0
\]
であり、核上四次のHesse係数は
\[
 \rho^2[xyz]\sum_i\{(\partial_i p_2)(Bx)_i+x_i(Bx)_i^2\}
 =-23\rho^2/2.
\]
高速補空間を最初の五座標とすると三次勾配は
$g=(12,-12,12,-6,6)^t$、$H_{0,\mathrm{fast}}\eta=-g$の解は
$\eta=(6,12,6,6,-6)^t$である。
停留消去の四次寄与は$g^t\eta/(2\rho)=-36/\rho$、従って
\[
 \kappa=-36/\rho-23\rho^2/2\ne0.
\]
ここではHesse規格化の寄与と停留消去の寄与を両方含めている。

第三境界成分には定理\ref{mh:four}の境界例の全$Q_i$を$\rho$倍して用いる。
核直線では$u=z(x_1+x_2),v=0$であり、制限式は二変数三次式と$x_3^3/3$の和なのでHesse係数は恒等的に零である。
高速三次勾配は$g=(0,0,0,6,0)^t$で、四次停留寄与は
$-g^tJ^{-1}g/(2\rho)=9/(4\rho)\ne0$である。
残る境界成分は添字の置換で得られる。
全例がFermat還元を持つので滑らかである。
小行列式チャート上で四次残余係数は正則関数で、その非消失は開であるから、各既約成分上で稠密となる。
\end{proof}

\section{cuspの正規化と導手}
\begin{proposition}\label{mh:cusp}
$A_2$模型の相対臨界値写像は
\[
 (z,t)\longmapsto(a,b,t)=(-3\eta z^2,2\eta z^3,t).
\]
像は$4a^3+27\eta b^2=0$であり、写像は有限正規化である。
$S=K[[t_1,\ldots,t_8]]$, $A=S[[z^2,z^3]]\subset B=S[[z]]$とすると
\[
 \mathfrak c_A=(a,b),\quad \mathfrak c_AB=(z^2),\quad
 \Omega_{B/A}=B/(z)\,dz,\quad I_9(d\pi)=(z).
\]
被約特異稜線の逆像は$(z^2)$、Jacobian特異スキームの逆像は$(z^3)$であり、
横断$\delta$不変量は一である。元の$\Delta$はこの近傍で単元を除いて$z$に一致する。
\end{proposition}
\begin{proof}
$h=y=0$と$\partial_z\sigma=0$から表示の写像を得る。
$b$に関するWeierstrass剰余$f(a)+bg(a)$の像は偶数次・奇数次が相殺しないので、核は表示の関係式だけである。
$B=A+Az$、$z=-3b/(2a)$なので有限双有理で、$B$は正則ゆえ正規化である。
$A$は$z$の一次係数が零の級数全体である。
$f\in A$について$fz\in A$の条件は定数項も零であることで、導手と$B/A=S\cdot z$が従う。
$d(z^2)=d(z^3)=0$は相対微分で$z\,dz=0$を与え、
$f\mapsto f'(0)$という$A$線形導分が逆を与える。
写像の微分の非自明列は$(-6\eta z,6\eta z^2)^t$である。
また$\operatorname{Hess}\sigma$の行列式は単元倍の$z$なので$\Delta$も同じイデアルを生成する。
$(a,b)B=(z^2)$、$(a^2,b)B=(z^3)$を代入すれば最後の主張を得る。
\end{proof}

\section{swallowtailの二つの稜線}
\begin{theorem}\label{mh:swallowtail}
$A_3$模型で$S_7=K[[t_1,\ldots,t_7]]$とする。
相対臨界スキームは$B=S_7[[a,z]]$であり、その標的写像は
\[
 b=-4\kappa z^3-2az,\qquad c=3\kappa z^4+az^2.
\]
その像$A$は次の既約被約判別式で定義され、$B$は有限正規化である：
\begin{align*}
 D_\kappa={}&256\kappa^3c^3-128\kappa^2a^2c^2+144\kappa^2ab^2c
 -27\kappa^2b^4+16\kappa a^4c-4\kappa a^3b^2.
\end{align*}
$R=a+6\kappa z^2$, $S=a+2\kappa z^2$, $p=SR^2$と置けば
\[
 \mathfrak c_B=(SR^2),\qquad \Omega_{B/A}=B/(R)\,dz,\qquad
 I_9(d\pi)=(R),\qquad J_{\mathrm{sing}}(A)B=(SR^3).
\]
導手の$A$内での生成元は$p,pz,pz^2$で、具体的には
\[
 p=a^3-4\kappa ac+\tfrac92\kappa b^2,\quad
 pz=-\tfrac b2(a^2+12\kappa c),\quad
 pz^2=\tfrac34ab^2-2a^2c+8\kappa c^2.
\]
元の$\Delta$は単元を除いて$R$に一致する。
\end{theorem}
\begin{proof}
臨界方程式を解くと表示の写像を得る。$z$は
$z^3+az/(2\kappa)+b/(4\kappa)=0$を満たすので、Weierstrass除法で$B$は標的上有限である。
重根を持つ四次多項式$\kappa z^4+az^2+bz+c$は全てこの表示を持つので判別集合は既約である。
$(a,b,c)=(0,-4\kappa,3\kappa)$では一つだけの二重根を持ち、$\partial_cD_\kappa\ne0$なので被約である。
重み$(2,3,4)$で準斉次な既約式なので形式的にも既約であり、像の核は$(D_\kappa)$である。
分数体では
\[
 z=-\frac{b(a^2+12\kappa c)}{2a^3-8\kappa ac+9\kappa b^2}.
\]
分母の引戻しは$2SR^2$で非零だから双有理である。
正則な$B$が正規化となる。

表示の三つの多項式恒等式から$pB\subset A$である。
逆に導手の元$q$について$q,zq\in A$であり、$a,t$を固定すると全$A$元の$z$微分は$(R)$に入る。
$q_z,(zq)_z\in(R)$より$q\in(R)$、$q=Rf$と書けば
$12\kappa zf+Rf_z\in(R)$より$f\in(R)$となる。
ここで$B/(R)$は整域である。従って$q\in(R^2)$。
$S=0$への$A$の像は$S_7[[z^2]]$であり、$q,zq$がともに偶数級数になるには$q\equiv0\pmod S$でなければならない。
$R,S$は相異なる素元なので$q\in(R^2)\cap(S)=(SR^2)$を得る。
$B=A+Az+Az^2$から導手の三生成元も従う。

$db=-2R\,dz$, $dc=2zR\,dz$で、$\partial_z\bmod R$が逆の導分を与えるため相対微分の表示を得る。
微分行列の二次小行列式は$-2R,2zR,-2z^2R$であり、自由な七方向を加えて$I_9=(R)$である。
Hessianは合同変換で移り、新しい形の行列式は単元倍の$R$なので$\Delta$も一致する。
直接代入で
\[
 (D_a,D_b,D_c)B=16\kappa SR^3(z^2,z,1)
\]
を得て、Jacobianイデアルの式が従う。
\end{proof}

\begin{corollary}
被約特異集合は二つの既約稜線
\[
 E_{\mathrm{cusp}}:(a^2+12\kappa c,\ 8a^3+27\kappa b^2),\qquad
 E_{\mathrm{double}}:(b,\ a^2-4\kappa c)
\]
の和であり、それぞれの逆像イデアルは$(R^2)$と$S(a,z)$である。
前者の正規化は$R=0$で、後者は滑らかで$S=0$から一般に二対一で覆われる。
二稜線の被約交わりは$a=b=c=0$、スキーム論的交わりは$S_7[[a]]/(a^2)$である。
\end{corollary}
\begin{proof}
前定理のJacobianイデアルから被約特異集合の記述が従う。
$R=0$では$(a,b,c)=(-6\kappa z^2,8\kappa z^3,-3\kappa z^4)$、
$S=0$では$(-2\kappa z^2,0,\kappa z^4)$である。
像の環はそれぞれ$S_7[[z^2,z^3]]$と$S_7[[z^2]]$なので核も表示の通りである。
さらに
\[
 a^2+12\kappa c=R^2,\quad8a^3+27\kappa b^2=4R^2(2a+3\kappa z^2),
\]
$b=-2zS$, $a^2-4\kappa c=S(a-6\kappa z^2)$を用いると引戻しの式を得る。
二イデアルの和は$(b,c,a^2)$である。
\end{proof}

\paragraph{範囲.}
以上は指定完備近傍の正規化と導手の計算であり、全Grassmannian上の臨界値像やその写像次数の分類ではない。
六つの自由な源変数との積を取った写像そのものは有限ではない。
$A_k$標準形は既知の特異点論であり、本稿の個別incidenceへの適用の文献上の初出は未確認である。

\section{四次残余の消失弧とA4}
\begin{theorem}\label{mh:a4}
滑らかな三次三様体の標識付き平面incidenceに、六次元$A_4$ファイバーの非空層がある。
$K=\CC((\tau))$, $\rho=\tau^2$とし、
\begin{align*}
 Q_1&=\rho\begin{pmatrix}(17-3q)/20&-3(q+1)/10\\-3(q+1)/10&(2-3q)/5\end{pmatrix},\\
 Q_2&=\rho\begin{pmatrix}(q+1)/4&(q-1)/2\\(q-1)/2&q\end{pmatrix},\qquad
 Q_3=\rho\begin{pmatrix}-3&1\\1&0\end{pmatrix}.
\end{align*}
と取る。$v_0^2=-952/1215$の一方を選び、
\[
 -405v(s)^2-952/3-190s v(s)+228s^2v(s)^2+144s^3v(s)^3=0,
 \quad v(0)=v_0,\quad q=5/3+\tau^3v(\tau^3)
\]
で形式弧を定める。この点の十次元標的近傍は、$h$を固定して
\[
 h^2=3\{q_5(y)+\kappa z^5+a_3z^3+a_2z^2+a_1z+a_0\},\qquad
 \kappa=\rho\frac{135(3q-1)}{2(4q-1)}\ne0
\]
となり、四つの係数は独立な標的座標である。
\end{theorem}
\begin{proof}
$H=\rho H_0(q)$の核は$k=(1,1,1,2,2,-1)$、最初の五次小行列式は
$D_q/10$、$D_q=(3q-1)(4q-1)$である。
$\Gamma(k)=0$であり、高速三次勾配は$g=(12,-12,12,-6,6)^t$。
$\eta=-H_{0,\mathrm{fast}}^{-1}g$を取ると
\begin{align*}
 B_4&=g^t\eta/2=-45(3q-5)^2/(4D_q),\\
 B_5&=\eta^t\Gamma''_{\mathrm{fast}}(k)\eta/2
 =-135(3q-5)^2(165q^2-54q-59)/(8D_q^2).
\end{align*}
これらは停留消去に由来する寄与である。Hesse規格化の寄与も計算する。
$L=x_1+x_2+2x_3,M=x_1+2x_2-x_3$、
$p_2=\tfrac12\sum x_i\{(Q_i/\rho)_{11}L^2+2(Q_i/\rho)_{12}LM+(Q_i/\rho)_{22}M^2\}$、
$p_3=L^3+M^3$とし、$N(p)_{ij}=-[x_i^2x_j]p$と置く。
$A=N(p_2)x,B=N(p_3)x$から
\begin{align*}
 C_4&=[xyz]\sum_i\{(\partial_i p_2)A_i+x_iA_i^2\}=(6q-17)/2,\\
 C_5&=[xyz]\sum_i\{(\partial_i p_2)B_i+(\partial_i p_3)A_i+2x_iA_iB_i\}
 =3(39q+175)/4.
\end{align*}
一般の核変数で定義した$C_4$の$(\eta,0)$方向微分は
$3(261q^3-2073q^2+605q+459)/(2D_q)$である。
従って$G_5=dC_4(k)[(\eta,0)]+C_5
=3(990q^3-2319q^2+24q+1093)/(4D_q)$。
停留解の三次補正の五次寄与は$(H_{0,\mathrm{fast}}\eta+g)^t\xi=0$なので、
残余四、五次は正確に
\[
 \kappa_4=\rho^{-1}B_4+\rho^2C_4,\qquad
 \kappa_5=\rho^{-2}B_5+\rho G_5
\]
となる。$\kappa_4=0$に上の弧を代入して$\tau^6$で割ると表示の$v$方程式を得る。
$v$微分は$-810v_0\ne0$なので形式陰関数定理が適用でき、$v=v_0-19s/81+O(s^2)$となる。
有理恒等式
\[
 B_4\left(G_5-\frac{135(3q-1)}{2(4q-1)}\right)-C_4B_5=0
\]
と$B_4=-\rho^3C_4$から$\kappa_5=\kappa=(810/17)\tau^2+O(\tau^5)\ne0$を得る。
全$Q_i$は$\tau^2$で割れるので三次式はFermat還元を持ち、一般ファイバーは幾何学的に滑らかである。
五変数Morse分裂と定数項一の単元の五乗根で中央関数を$\kappa z^5$にする。
残余Jacobian商$K[[z]]/(z^4)$において、標的方向$xyz,ux_2x_3,x_1uv,u^3$の先頭類は
$1,z,3z^2,12z^3$であり、行列式は$36$である。
軌道方向は平面移動で相殺されるので十次元sliceに降りる。
定理\ref{mh:versal}と同じ小拡大帰納で四係数を吸収して正規形を得る。
\end{proof}

\begin{corollary}
上の二パラメータ族では$\rho D_q\ne0,\kappa_4=0$なら常に$\kappa_5\ne0$であり、$A_5$点はない。
これは九係数族全体の$A_5$不存在を意味しない。
$A_4$近傍の臨界値写像は
\[
 a_1=-5\kappa z^4-3a_3z^2-2a_2z,\qquad
 a_0=4\kappa z^5+2a_3z^3+a_2z^2
\]
であり、五次多項式の判別式で定義される像の有限正規化である。
三次以下の残余が消える層は$a_3=-10\kappa z^2,a_2=20\kappa z^3$。
この七次元空間では$z=0$が$A_4$層、$z\ne0$は$A_3$層である。
\end{corollary}
\begin{proof}
$C_4=0$の$q=17/6$では$B_4\ne0$なので、四次消失上で$C_4$は非零である。
前証明の恒等式から五次は表示の非零式となる。
臨界値写像の式は$g=g'=0$を解けばよい。
$z$はmonic四次関係を満たすので有限であり、一般に一つの二重根を持つ多項式の重根指定なので双有理である。
判別式は既約被約で重み$(2,3,4,5)$の準斉次式であり、形式的にも既約である。
被約性は例えば$(a_3,a_2,a_1,a_0)=(0,0,-5\kappa,4\kappa)$で
$\partial_{a_0}\operatorname{Disc}=800000\kappa^7\ne0$からも確認できる。
$g''=g'''=0$を解いて層の式を得る。
そこで$g''''=120\kappa z,g'''''=120\kappa$であるから最後の主張が従う。
\end{proof}

\section{D4点、二次錐、主成分の非正規性}
\begin{theorem}\label{mh:d4}
$K=\CC((\rho))$上で$Q_1=\rho I,Q_2=\rho\operatorname{diag}(1,-1),
Q_3=\rho\left(\begin{smallmatrix}0&1\\1&0\end{smallmatrix}\right)$とすると、
三次三様体は滑らかで、平面incidenceは$D_4$ファイバーを持つ。
$h$を固定した標的・相対形式座標で
\[
 h^2=3\{\rho(y_1y_4+y_2y_3)+6st(s-t)+\alpha st+\beta s+\gamma t+\delta\}
\]
となり、四係数は独立な標的座標で、他に六つの自由な標的座標がある。
臨界源は$K[[\alpha,s,t,\tau_1,\ldots,\tau_6]]$で、写像は
\begin{align*}
 \beta&=-12st+6t^2-\alpha t,&
 \gamma&=-6s^2+12st-\alpha s,&
 \delta&=12st(s-t)+\alpha st.
\end{align*}
この近傍で$(\Delta)=I_9(d\pi)=(R)$、
$R=(\alpha+12(s-t))^2+144st$である。
従って六次退化因子は通常二次曲面錐と六次元円板の積で、特異スキームは$(\alpha,s,t)$である。
\end{theorem}
\begin{proof}
$H$の核は$n(s,t)=(s,t,-s,t,-t,-s)$で、最初の四座標上のHessianは
$\rho\left(\begin{smallmatrix}0&0&0&1\\0&0&1&0\\0&1&0&0\\1&0&0&0\end{smallmatrix}\right)$、行列式$\rho^4$である。
核上三次は$f_3=6st(s-t)$で、判別式$1296$を持つ。
Jacobian商は$K\langle1,s,t,st\rangle$、$s^2=t^2=2st$、$(s,t)^3=0$である。
従って四次以上の各誤差は、次数を二つ下げた係数を持つ二偏微分の線形結合となる。
Morse分裂後にこの誤差を次数順に除けば中央関数は$f_3$へ右同値となる。
$u=s-t,v=s+t$で$f_3=3u(v^2-u^2)/2$だから幾何学的に$D_4$型である。
Fermat還元から周囲三様体も滑らかである。

高速停留解の二次項は$(6st/\rho)(1,1,1,-1)$である。
核上の$L=sx_1-sx_2-tx_3,M=tx_1+tx_2-sx_3$に対し
$p_2=x_1(L^2+M^2)/2+x_2(L^2-M^2)/2+x_3LM$と置く。
二次規格化行列$B_2$の跡は$-(s^2+st)$である。
標的方向$xyz,ux_2x_3,vx_2x_3,x_2uv$の残余商での列は、基底$(1,s,t,st)$について
\[
 \begin{pmatrix}1&0&0&0\\0&1&0&0\\0&0&1&0\\-3\rho&6/\rho&6/\rho&-4\end{pmatrix}
\]
となり、行列式$-4$で全射である。
軌道方向を除いた十次元sliceにも降りるので、小拡大で誤差を
$e_0+e_1s+e_2t+e_3st+A\partial_sf_3+B\partial_tf_3$へ分ける帰納が正規形を与える。
全過程で$h$を変えず、方程式の単元倍を用いない。

$g=g_s=g_t=0$を解くと表示の臨界値写像を得る。
残余Hessianは
$\left(\begin{smallmatrix}12t&12(s-t)+\alpha\\12(s-t)+\alpha&-12s\end{smallmatrix}\right)$で行列式$-R$。
$d\delta=-s\,d\beta-t\,d\gamma-st\,d\alpha$から微分の最大小行列式も$R$で生成される。
四つの高速方向の行列式は単元なので元の$\Delta$も同じイデアルを与える。
$A=\alpha+12(s-t)$と置くと$R=A^2+144st$であり、Jacobianイデアルは$(A,s,t)$である。
\end{proof}

\begin{proposition}
この近傍の$A_3$層の被約閉包は三直線
\[
 s=0,\ \alpha=12t;\qquad t=0,\ \alpha=-12s;\qquad
 s=-t,\ \alpha=12t
\]
と六次元円板の積である。頂点以外は全て$A_3$型で$A_4$点を持たない。
三枝は全て大域的主成分$M$に属し、三つの境界成分ではない。
$S=K[[\tau_1,\ldots,\tau_6]]$とすると主成分の被約完備局所環は
\[
 A_M=S[[X,Y,Z]]/(XY,XZ,YZ),\quad
 X=s+\alpha/12,\ Y=t-\alpha/12,\ Z=t-s-\alpha/12.
\]
正規化は$\bar A_M=S[[X]]\times S[[Y]]\times S[[Z]]$、導手は$(X,Y,Z)$、
$\bar A_M/A_M\simeq S^2$である。従って$M$はこの$D_4$層に沿って非正規で、横断$\delta$不変量は二である。
\end{proposition}
\begin{proof}
二次錐の頂点外で残余Hessianの階数は一である。
核上$f_3$の消失は核が$(0,1),(1,0),(1,1)$のいずれかとなることなので三直線を得る。
第一の$(s,t,\alpha)=(0,u,12u)$では平行移動後の関数は
$6(u+z)x^2-6xz^2$であり、停留消去後$-3z^4/(2(u+z))$となる。
第二では$3z^4/(2(u+z))$、第三では$-3z^4/(2(u+z))$となるから全て四次残余非零である。
元の$Q_i$は全て可逆なので、$\det Q_i=0$に含まれる境界成分はこの近傍を通らない。
表示の可逆線形変換は三直線を座標三軸へ移す。
有限型の被約環の完備化はexcellent性により被約なので、これが主成分の完備局所環である。
積環への像は三つの定数項が一致する組であり、成分冪等元の添加は整である。
各成分が正規なので表示の積環が有限正規化となる。
二つの定数項の差が正規化商を$S^2$へ同定し、冪等元を掛けても像に戻る条件は共通定数項が零であることなので導手を得る。
\end{proof}

\section{相対二重被覆と古典臨界値}
$A=\CC[[s,P,Q,T]]$、$z=(\lambda,\mu)$とし、平面incidenceを
\[
 B=A[[z,h]]/(h^2-3\sigma(s,P,Q,T,z))
\]
と書く。$\partial_s\sigma(0)=1$なので$B$は正則な十六次元形式環である。
$A[[z]]$上は基底$1,h$の有限自由二重被覆であるが、
$A$上の相対臨界スキームはその分岐点スキーム全体とは異なる。
\begin{proposition}
$A$上の相対臨界イデアルは$(h,\partial_z\sigma)$であり、
その環は$\CC[[T,z]]$という正則九次元環となる。
$\sigma=0$を$s=S_{\rm cl}(P,Q,T,z)$と解けば、この臨界環の別表示は
\[
 \CC[[P,Q,T,z]]/(\partial_zS_{\rm cl}),\qquad h=0,\quad s=S_{\rm cl}.
\]
元の源側の六座標を保持して二重被覆を$h^2=\Phi_A(s,r,a,c)$と表示した場合も、
古典零点$s=S_{\rm cl}(r,a,c)$が同じなら、一意な単元$\upsilon$により
$\Phi_A=\upsilon(s-S_{\rm cl})$となる。
従って分岐零スキームと臨界値写像は古典族で決まり、未知量子級数に依存しない。
\end{proposition}
\begin{proof}
連続相対微分は七生成元$dh,dz$と一関係を持つので、第六Fittingイデアルが表示式となる。
$h=0,\sigma=0$から$s$を解き、さらに六つの微分を$P,Q$について解く。
原点でそのJacobianは恒等行列だから臨界環は正則である。
最後の因数分解は$\Phi_A(S_{\rm cl},r,a,c)=0$と$\partial_s\Phi_A(0)\ne0$から従う。
基底$1,h$のtrace pairingは$\operatorname{diag}(2,2\Phi_A)$なので被覆の判別イデアルは$(\Phi_A)$であり、
相対臨界イデアルを$h=0$で計算すると$\partial_z\Phi_A=-\upsilon\partial_zS_{\rm cl}$となる。
\end{proof}

\begin{proposition}[一般T方向のMorse模型]
$K=\operatorname{Frac}\CC[[T_1,T_2,T_3]]$へ移り、$s=P=Q=0,z=0$で完備化する。
$M(T)=\left(\begin{smallmatrix}0&T_3&T_2\\T_3&0&T_1\\T_2&T_1&0\end{smallmatrix}\right)$とすると
\[
 \sigma_{zz}=\begin{pmatrix}0&M(T)\\M(T)&0\end{pmatrix},\quad\det M=2T_1T_2T_3.
\]
一意な停留点とその値$d$は
\[
 \lambda_*=-M^{-1}Q+O((s,P,Q)^2),\quad
 \mu_*=-M^{-1}P+O((s,P,Q)^2),\quad
 d=s-P^tM^{-1}Q+O((s,P,Q)^3)
\]
を満たす。$h$を固定する相対形式座標で
$h^2=3(\sum_iU_iV_i+d)$となる。
臨界値因子$d=0$は滑らか被約で、臨界値写像はこの局所因子への同型である。
\end{proposition}
\begin{proof}
制限式は$D/3+\ell^3+m^3+\ell m\sum T_ix_i$であり、平面二次項が正確に$\lambda^tM\mu$である。
Hesse規格化の効果は平面四次以上なのでHessianに高い$T$補正はない。
逆関数定理で停留点を解き、二次項へ代入すると表示の値を得る。
$\partial_sd(0)=1$であり、相対Morse帰納が残余の非退化形式を三双曲対へ移す。
\end{proof}

\section{二つの退化因子の交点}
\begin{theorem}
$K_1=\operatorname{Frac}\CC[[T_1]]$, $R_1=K_1[[T_2,T_3]]$上の$T$軸族には、
$T,h$を保つ形式座標で
\[
 \sigma=U_1V_1+U_2V_2+T_2T_3AB
\]
という正規形がある。相対特異スキームの環は
$K_1[[T_2,T_3,A,B]]/(T_2T_3A,T_2T_3B)$であり、
被約な三成分$(T_2),(T_3),(A,B)$を持つ。
全空間は八次元の正規完備局所環で、その特異スキームは
$h=U=V=0$と、四変数$T_2,T_3,A,B$の任意の三つの積が零という被約な六座標平面の和である。
\end{theorem}
\begin{proof}
$a=\lambda_1,b=\mu_1$とし、他の四変数$w$のHessianは行列式$T_1^4$を持つ。
$\partial_w\sigma=0$を解いてMorse分裂すると、$\sigma=U_1V_1+U_2V_2+g(a,b)$となる。
$T_3=0,\lambda_3=\mu_3=0$上の制限式は二元三次式と$x_3^3/3$の和である。
二元三次式を近恒等$\GL_2$でFermatへ戻せるのでHesse係数は恒等的に零である。
残る二法微分を$\lambda_2,\mu_2$について解くと、そのJacobianは$T_1I$であり、二次元臨界曲面を得る。
停留解の一意性から$g\bmod T_3=0$である。添字交換で$g\bmod T_2=0$も得る。
従って$g=T_2T_3q$。
二次停留解は$\lambda_2=-T_2a/T_1,\lambda_3=-T_3a/T_1$および同じ$\mu$式なので、
$g_2=-2T_2T_3ab/T_1$となる。
$q$のHessianは可逆で、二変数Morse補題により$q=AB$とできる。

相対Jacobianと全Jacobianを取れば表示のスキームを得る。
単項式イデアルの等式
\[
 (T_2T_3A,T_2T_3B)=(T_2)\cap(T_3)\cap(A,B),\qquad
 (\text{四変数の三重積})=\bigcap_{i<j}(z_i,z_j)
\]
が被約性を与える。全空間は超曲面でCohen--Macaulay、特異集合の余次元は六なので正規である。
\end{proof}

\begin{proposition}[退化因子一般点の二次jet]
$A=T_1,B=T_2$を係数体$K=\operatorname{Frac}\CC[[T_1,T_2]]$の単元とし、$\tau=T_3$と置く。
標的イデアル$\mathfrak b=(s,P,Q,\tau)$に対し、残る二平面変数を$\lambda_1=A\alpha,\mu_1=A\beta$に選ぶと、
四高速方向の消去後の関数は
\begin{align*}
 g\equiv{}&s+(AP_1-BP_2)\alpha+(AQ_1-BQ_2)\beta
 +6BQ_2\alpha^2+6BP_2\beta^2\\
 &-\{2AB\tau+(A^3+B^3)s\}\alpha\beta
 \pmod{\mathfrak b^2+\mathfrak b(\alpha,\beta)^3}.
\end{align*}
$P_3,Q_3$方向の一次変形関数は全次数で零である。
二変数二次jetへのKodaira--Spencer写像は全射で、列$s,P_1,P_2,Q_1,Q_2,\tau$の行列式は$72A^3B^3$である。
しかし全変形の接空間は$K[[\alpha,\beta]]$なので、この有限次元標的は全形式変形に対してversalではない。
\end{proposition}
\begin{proof}
中央の臨界曲面は$\lambda_3=\mu_3=0$、
$\lambda_2=-B\alpha+6B\beta^2+O(3)$、
$\mu_2=-B\beta+6B\alpha^2+O(3)$である。
制限二元三次式の二次摂動は
$\alpha\beta(Ax_1+Bx_2)(Ax_1-Bx_2)^2$なので二元規格化行列は
\[
 I+\alpha\beta\begin{pmatrix}-A^3&A^2B\\AB^2&-B^3\end{pmatrix}+O(3).
\]
摂動をこの行列で変換した後の$xyz$係数を取ると表示の八方向の変形関数を得る。
$P_3,Q_3$摂動は二元三次式のままで$x_3$を生じないから全次数で消える。
行列式は直接の係数比較による。
中央ファイバーの方程式$h^2-3(U_1V_1+U_2V_2)$のJacobian商は$K[[\alpha,\beta]]$である。
\end{proof}

\paragraph{微分の階数.}
$T$軸の原点平面section上で、普遍臨界値写像の微分は
$(\delta T,\delta\lambda,\delta\mu)\mapsto
(0,-M(T)\delta\mu,-M(T)\delta\lambda,\delta T)$である。
従って階数は$T_1T_2T_3\ne0$なら九、$T\ne0,T_1T_2T_3=0$なら七、$T=0$なら三である。
ここで$M$の非零二次小行列式$-T_i^2$により、可能階数が$3,2,0$であることを用いた。

\begin{corollary}[忠実平坦性とArtin点の持上げの相違]
平面incidenceの環準同型$A\to B$は忠実平坦で、閉ファイバーは六次元である。
六つの自由な源変数を加えた場合も忠実平坦で、閉ファイバーは十二次元である。
どちらも原点で形式的に滑らかではなく、同じArtin基底上の点の写像は一般に全射でない。
\end{corollary}
\begin{proof}
$A[[z]]$は$A$上忠実平坦で、$B$はその上基底$1,h$の有限自由加群だから主張を得る。
閉ファイバーは七変数正則環の非零一式による商なので六次元である。
原点で$ds$方向は接写像の像に入らない。
実際$A\to\CC[\epsilon]/(\epsilon^2)$を$s\mapsto\epsilon,r\mapsto0$とすると、
$h,z$の全像が$(\epsilon)$に入る一方、$3\sigma=3s+O(2)$なので関係式を満たせない。
自由変数の追加はこの障害を消さない。忠実平坦性は同一基底上の点の全射性を意味しない。
\end{proof}

\section{固定源座標の分岐jetと形式型}
本節は全混合可換代表の次の入力を明示して用いる。
対角平方線形項を除き、$e_i=p_iu+q_iv$、$m_i=a_iu^2+d_iuv+c_iv^2$とする。
$h$の偶級数である純曲率代表が、$F_0=D/3+u^3+v^3$に
\[
 \frac{h^2}3xyz+\sum_i e_ix_jx_k+\sum_i m_ix_i
 -\frac{h^2}{12}\sum_i e_i^2x_i+\frac{h^2}4e_1e_2e_3
\]
を加えた式に$h^3$を法として一致すると仮定する。
さらに全次数の代表の追加係数が
\begin{gather*}
 b-1,\gamma-1=O(h^6),\qquad U_e=-h^2/12+O(h^8),\\
 S_e=h^2/4+O(h^8),\qquad V_e,T_e=O(h^4),\qquad R_e=O(h^6).
\end{gather*}
という選択則を持つものとする。
これらは形式量子化からの比較の入力であって、通常の三次式の軌道計算だけからは従わない。
現行の\href{https://github.com/uedakazushi/nc-cubic2/blob/main/papers/sklyanin_mixed_formal_comparison.tex}{混合形式比較稿}で、有限gauge・相対window・全六級数の定理として証明する。
以下では、この入力を持つ任意の代表に対する幾何学的帰結を証明する。

\begin{theorem}\label{mh:fibre}
canonical Fermat sliceの九混合座標を零にして$(p,q,d)$を消去すると、一意な分岐式
$h^2=\Phi(a,c)$を得る。その三、六、九次項は
\begin{align*}
 \Phi&=\frac23(P_a+P_c)+\frac2{27}R_6+\Phi_9+O((a,c)^{12}),\\
 \Phi_9&=-\frac2{243}(P_aE_a+P_cE_c)-\frac4{243}(P_aE_c+P_cE_a)
       +\frac4{243}(P_a+P_c)J_6,
\end{align*}
ただし
\[
 P_a=\prod_i a_i,\quad E_a=\sum_{i<j}a_i^3a_j^3,\quad
 J_6=\sum_{i\ne j}a_i^2c_i a_jc_j^2,\quad
 R_6=\sum_i a_i^2c_i^2(a_jc_k+a_kc_j)
\]
で、$P_c,E_c$も同様である。
有限対角群の指標条件を持つ場合、$h$を固定した接写像恒等の形式座標で
\[
 \CC[[h,A,C]]/\bigl(h^2-\tfrac23(A_1A_2A_3+C_1C_2C_3)\bigr)
\]
となる。元の源係数での六、九次項は一般に零ではない。
\end{theorem}
\begin{proof}
原点で九混合係数から九slice座標への微分は恒等なので、形式逆関数定理で一意に消去できる。
$h$の偶性から残る式は$t=h^2$の級数で、$\partial_tS(0)=1/3$である。
従って$S=\text{単元}\,(t-\Phi)$と一意に書ける。
$u\mapsto u-A/3,v\mapsto v-C/3$、$A=\sum a_ix_i,C=\sum c_ix_i$を用いると
$p_i=2a_ja_k/3+O(3),q_i=2c_jc_k/3+O(3),d_i=O(3)$、
$S=t/3-2(P_a+P_c)/9+O(4)$を得る。

古典軌道の計算を源係数を固定して行う。
$z=(x_1,x_2,x_3,u,v)$、$\kappa=(1/3,1/3,1/3,1,1)$とし、
$F_0(Mz)$の二十五の非squarefree係数を
$H=F_0+u^2\sum a_ix_i+v^2\sum c_ix_i$に一致させる。
$M\equiv I$は一意で、次数$n$の増分は
\[
 (\delta M_{ij})_n=-\frac1{3\kappa_i}[z_i^2z_j](F_0(M_{<n}z)-H)_n.
\]
この再帰は係数パラメータで割らない。古典分岐式は
$\Phi^{\rm cl}=18\sum_{r=1}^5\kappa_rM_{r1}M_{r2}M_{r3}$である。
必要な行列成分は
\begin{align*}
 M_{ii}&=1-(a_i^3+c_i^3)/27+O(6),\\
 M_{ij}&=-(a_i^2a_j+c_i^2c_j)/9+N_{ij}+O(8)\quad(i\ne j),\\
 N_{ij}&=\{4(a_i^5a_j+c_i^5c_j)-2(a_i^2a_jc_i^3+a_i^3c_i^2c_j)
                    -6a_ia_j^2c_ic_j^2\}/243,\\
 M_{ui}&=a_i/3-a_i^4/27+
 \{7a_i^7+a_i\sum_{j\ne i}a_j^6-3a_i^4c_i^3+3c_i\sum_{j\ne i}a_j^4c_j^2\}/729+O(8),
\end{align*}
$v$行は$a,c$を交換した式である。
表示の再帰を次数七まで代入して得られ、全二十五方程式の残差は七次以下で零である。
$x$行の$xyz$係数は次数$3+3+3$または$0+3+6$、
$u,v$行では$1+1+7$または$1+4+4$から生じる。
八次以上の未知行列成分はそれぞれ次数十一、十以上にしか寄与しない。
代入して
\[
 \Phi^{\rm cl}=\frac23(P_a+P_c)+\frac2{27}R_6+
 \frac2{243}(P_aE_a+P_cE_c)+\frac4{243}(P_a+P_c)J_6+O(10)
\]
を得る。この有限多項式計算は、入力行列を使わない現行検証器でも再生成する。

実際の代表で$\operatorname{ord}t\ge3,\operatorname{ord}(p,q)\ge2,\operatorname{ord}d\ge3$を用いる。
純$u,v$三次補正は九次以上なので、一意な$\GL_2$正規化$I+O(9)$で除ける。
それが$a,c$へ与える変化は十次以上、$p,q$へは十一次以上である。
九次分岐式へ影響する量子項は$-t\sum e_i^2x_i/12$だけであり、有効係数は
\[
 \widetilde a_i=a_i-tp_i^2/12+O(10),\quad
 \widetilde c_i=c_i-tq_i^2/12+O(10),\quad
 \widetilde d_i=d_i-tp_iq_i/6+O(10).
\]
従って$t=\Phi^{\rm cl}(\widetilde a,\widetilde c)$である。
七次係数増分は$\delta a_i=-t_3a_j^2a_k^2/27$、同じ$c$式となるので
三次項のTaylor一次効果は$-4(P_a+P_c)(E_a+E_c)/243$である。
これを古典九次項へ加えると表示式を得る。

対角群$\lambda_i^3=\eta^3=\theta^3=1$、$\Lambda=\prod\lambda_i$に対して
$a_i\mapsto\lambda_i^{-1}\eta a_i,c_i\mapsto\lambda_i^{-1}\theta c_i$、
$\Phi\mapsto\Lambda^{-1}\Phi$という指標を持つとする。
単項式$a^Ac^C$は$A_i+C_i\equiv1\pmod3$、$\sum A_i\equiv\sum C_i\equiv0\pmod3$を満たす。
従って全次数は三の倍数で、四、五、七、八、十、十一次はない。
また全ての添字がsupportに現れ、二群への配分の鳩ノ巣原理から
各単項式は$a_ia_j$または$c_ic_j$で割れる。
これは$f_3=2(P_a+P_c)/3$のJacobianイデアルである。
次数$N$の誤差を$V_N(f_3)$と書き、有限群で平均した不変な次数$N-2$ベクトル場を選ぶ。
$z\mapsto z-V_N(z)$で誤差を除き、次数順に反復する。
最初の修正は四次であり、各有限jetは有限回で安定する。
これが$h$を固定する同変形式右同値を与える。
\end{proof}

\begin{corollary}
最初の座標修正は
$\delta a_k=-C_k\sum_{i\ne k}A_iC_i^2/18$、
$\delta c_k=-A_k\sum_{i\ne k}C_iA_i^2/18$と取れ、六次項を消す。
形式ファイバーは正規六次元整域で、接空間次元七、重複度二。
特異スキームは
\[
 (h,A_1A_2,A_1A_3,A_2A_3,C_1C_2,C_1C_3,C_2C_3)
\]
という被約な九座標二次元平面の和であり、一般横断型は四次元通常二重点である。
\end{corollary}
\begin{proof}
$df_3(\delta a,\delta c)=-2R_6/27$を展開すれば修正の式を得る。
$f_3$は位数三なので正規環$\CC[[A,C]]$の分数体でも平方でなく、二重被覆は整域である。
偏微分の表示イデアルは九平面の素イデアルの共通部分であり、特異集合の余次元は四である。
超曲面の$(S_2)$と$(R_1)$から正規性を得る。
一般点では$A_i,C_j$が単元なので、横断方程式は$h$と残る四変数の非退化二次形式である。
\end{proof}

\section{一混合Hesse三次式の完全resultant}
この節は形式比較から独立した三次式の幾何である。
\[
 F_{c,\kappa,\alpha}=(x^3+y^3+z^3)/3+cxyz+
 \kappa(\alpha uvx+u^3+v^3),\qquad d=1+c^3
\]
と置く。gradient resultantは$\operatorname{Res}(x^2,y^2,z^2,u^2,v^2)=1$で正規化する。
\begin{theorem}
全パラメータで多項式恒等式
\[
 \operatorname{Res}(F_x,F_y,F_z,F_u,F_v)
 =3^8\kappa^{32}d^3(\kappa\alpha^3+9)^3(\kappa\alpha^3+9d)^9
\]
が成立する。
$c\kappa d\ne0$なら特異条件は$(\kappa\alpha^3+9)(\kappa\alpha^3+9d)=0$であり、
前者は三つ、後者は九つの通常二重点を与える。
$c=0,\kappa\ne0$では二因子が一致して重複度十二となり、三特異点の局所型は
$p^2+q^2+y^3+z^3$である。
\end{theorem}
\begin{proof}
$\kappa d\ne0$の開集合でまず計算する。
$x=0$上では$F_y=F_z=0$から$y=z=0$、$F_u=F_v=0$から$u=v=0$なので射影的特異点はない。
$x=1$の勾配式は
\[
 y^2+cz=z^2+cy=0,\quad u^2+\alpha v/3=v^2+\alpha u/3=0,\quad
 1+cyz+\kappa\alpha uv=0.
\]
最初の対は$(0,0)$または$(-c\omega,-c\omega^2)$、$\omega^3=1$、
次の対は$(0,0)$または$(-\alpha\eta/3,-\alpha\eta^2/3)$、$\eta^3=1$である。
従って最後の式は二種類の因子を与え、点数はそれぞれ三、九となる。
四変数Hessianは二ブロック
\[
 \begin{pmatrix}2y&c\\c&2z\end{pmatrix},\qquad
 \kappa\begin{pmatrix}6u&\alpha\\\alpha&6v\end{pmatrix}
\]
に分かれ、その行列式はそれぞれ全体として$-3c^2\kappa^2\alpha^2$、
$9c^2\kappa^2\alpha^2$である。よって全点が通常二重点である。

定数と重複度をnormで決める。
$K[r,s]/(r^2+ts,s^2+tr)$の基底$(1,r,s,rs)$で$M_t=$乗法$\,rs$の行列は
\[
 M_t=\begin{pmatrix}0&0&0&0\\0&t^2&0&0\\0&0&t^2&0\\1&0&0&t^2\end{pmatrix}.
\]
従って上の四二次式による十六次元商上で$F_x$のnormは
\[
 \det(I_{16}+cM_c\otimes I_4+\kappa\alpha I_4\otimes M_{\alpha/3})
 =d^3(1+\kappa\alpha^3/9)^3(d+\kappa\alpha^3/9)^9.
\]
無限遠の先頭四式のresultantは$(3\kappa)^{16}$であり、
Poisson積公式\cite{Resultant}で次数二の$F_x$を加えるとその平方が掛かる。
分母$9^{12}$を整理して$3^8$を得る。
稠密開集合の計算は全パラメータの多項式恒等式へ延長する。
$c=0$では三つの$(u,v)$点でそのHessianは非退化で、$y,z$は独立三乗項のままなので表示の局所型となる。
$\kappa=0$なら頂点$\PP^1$を持つcone、$d=0$なら$u=v=0$上に三つの特異点が残ることも式から分かる。
\end{proof}

\begin{corollary}
$c=h^2/3,d=1+h^6/27,\kappa=d$とした候補式は
\[
 D/3+h^2xyz/3+d(\alpha uvx+u^3+v^3)
\]
であり、そのresultantは
$3^8d^{44}(d\alpha^3+9)^3(\alpha^3+9)^9$である。
$d\ne0$なら特異条件は
$(\alpha^3+9)((27+h^6)\alpha^3+243)=0$となる。
形式的には$b(h)^3=d,b(0)=1$により$u,v$を$b^{-1}$倍すれば
$D/3+h^2xyz/3+u^3+v^3+\alpha b(h)uvx$に移る。
この座標恒等式だけから有限$h$での圏同値や比較核の代数化は従わない。
\end{corollary}

\paragraph{有限係数候補の強制される規格化.}
$c=1/3$では$d=28/27$である。
この狭い形の候補で二つの悪い$\alpha^3$値を$-9,-243/28$へ一致させるには、
一般式の比が$d$であることから$\kappa=28/27$が強制される。
これは候補式のパラメータ照合であり、有限$h$での圏比較の存在証明ではない。

\begin{thebibliography}{9}
\bibitem{Resultant} C.~D'Andrea, T.~Krick and A.~Szanto,
\emph{Multivariate subresultants in roots}, J. Algebra 302 (2006), 16--36,
\url{https://arxiv.org/abs/math/0501281}, equation (2).
\bibitem{Stacks} The Stacks Project, Tags 031S (Serre's criterion), 07QS (excellent rings),
\url{https://stacks.math.columbia.edu}.
\bibitem{Arnold} V.~I.~Arnol'd, S.~M.~Gusein-Zade and A.~N.~Varchenko,
\emph{Singularities of Differentiable Maps}, vol.~I, Birkh\"auser, 1985.
\end{thebibliography}
\end{document}
